\documentclass[12pt,a4paper]{article}

\usepackage[a4paper,margin=1in]{geometry}
\usepackage{amsmath,amssymb}
\usepackage{zed-csp}
\usepackage{longtable}
\usepackage{titlesec}
\usepackage{setspace}

\setstretch{1.2}

\title{
\textbf{NeuroScript}\\
AI-Based Handwriting Analysis for Neurological Detection
}

\author{
INTERNATIONAL ISLAMIC UNIVERSITY, ISLAMABAD\\
Department of Software Engineering
}

\date{}

\begin{document}

\maketitle

\begin{center}
\textbf{COURSE:} Formal Methods\\
\textbf{INSTRUCTOR:} Ms. Saba Taimouri
\end{center}

\vspace{1cm}

\textbf{Submitted By:}
\begin{itemize}
    \item Areesha Abbasi (4508-FOC/BSSE-F22-A)
    \item Sajjal Khalil (4488-FOC/BSSE-F22-A)
    \item Hajira Gul (4454-FOC/BSSE-F22-A)
    \item Munazza Kanwal (4475-FOC/BSSE-F22-A)
    \item Noor ul Harmain (4204-FBAS/BSSE-F21-A)
\end{itemize}

\newpage

\tableofcontents

\newpage

\section{Introduction}

NeuroScript is an AI-powered web-based system that performs non-invasive early screening for neurological conditions — specifically OCD and ADHD — by analysing handwriting samples submitted by clinicians, students, or researchers.

The system accepts real-time stylus/finger input or digitised paper samples, extracts neuromotor features (stroke speed, spacing, letter formation), runs an LSTM-based classification model, and produces an explainable prediction via SHAP/LIME.

This document presents a complete formal Z specification covering:

\begin{itemize}
    \item Basic types and free-type enumerations
    \item State-space schema
    \item Initial-state schema
    \item Operational schemas for all major functions
    \item Error-reporting schemas and robust combined operations
\end{itemize}

\section{Basic Types and Free-Type Declarations}

\subsection{Basic Sets}

The following application-independent sets are introduced as basic types:

\[
[USER, SAMPLEID, FEATUREVECTOR, EXPLANATION]
\]

\begin{itemize}
    \item USER — The set of all registered system users.
    \item SAMPLEID — Unique identifiers for handwriting samples.
    \item FEATUREVECTOR — Extracted neuromotor feature vectors.
    \item EXPLANATION — SHAP/LIME explanation objects.
\end{itemize}

\subsection{Free-Type Enumerations}

\begin{zed}
ROLE ::= Clinician | Researcher | Student | Admin
\end{zed}

\begin{zed}
CONDITION ::= Normal | PotentialADHD | PotentialOCD | Inconclusive
\end{zed}

\begin{zed}
SAMPLESTATUS ::= Pending | Processing | Analysed | Archived
\end{zed}

\begin{zed}
REPORT ::= success | alreadyExists | notFound | invalidInput \\
\t1 | unauthorised | limitExceeded | processingError
\end{zed}

\subsection{Global Constants}

\[
MAX\_SAMPLES\_PER\_USER == 500
\]

\[
MIN\_STROKE\_POINTS == 50
\]

\section{State-Space Schema}

The central state of NeuroScript tracks registered users, uploaded samples, extracted features, model predictions, and generated explanations.

\begin{schema}{NeuroScript}
registeredUsers : \power USER \\
userRoles : USER \pfun ROLE \\
samples : SAMPLEID \pfun USER \\
sampleStatus : SAMPLEID \pfun SAMPLESTATUS \\
features : SAMPLEID \pfun FEATUREVECTOR \\
predictions : SAMPLEID \pfun CONDITION \\
explanations : SAMPLEID \pfun EXPLANATION \\
sampleCount : USER \fun \nat
\where
dom userRoles = registeredUsers \\
dom samples \subseteq dom sampleStatus \\
dom predictions \subseteq dom features \\
dom explanations \subseteq dom predictions \\
\forall u : registeredUsers @ \\
\t1 sampleCount~u = \# \{ s : dom samples | samples~s = u \} \\
\forall u : registeredUsers @ \\
\t1 sampleCount~u \leq MAX\_SAMPLES\_PER\_USER
\end{schema}

\subsection*{Invariant Explanation}

\begin{itemize}
    \item Every user with a role is a registered user and vice versa.
    \item A sample can only have a status if it belongs to the samples map.
    \item Features may only exist for known samples.
    \item Predictions require features.
    \item Explanations require predictions.
    \item Sample count must not exceed the per-user limit.
\end{itemize}

\section{Initial-State Schema}

When the system is first deployed, no users, samples, or analysis results exist.

\begin{schema}{InitNeuroScript}
NeuroScript
\where
registeredUsers = \emptyset \\
userRoles = \emptyset \\
samples = \emptyset \\
sampleStatus = \emptyset \\
features = \emptyset \\
predictions = \emptyset \\
explanations = \emptyset \\
sampleCount = \emptyset
\end{schema}

\section{Operational Schemas}

\subsection{User Registration}

\subsubsection*{Normal Case}

\begin{schema}{RegisterUser}
\Delta NeuroScript \\
newUser? : USER \\
role? : ROLE \\
result! : REPORT
\where
newUser? \notin registeredUsers \\
registeredUsers' = registeredUsers \cup \{newUser?\} \\
userRoles' = userRoles \cup \{newUser? \mapsto role?\} \\
samples' = samples \\
sampleStatus' = sampleStatus \\
features' = features \\
predictions' = predictions \\
explanations' = explanations \\
sampleCount' = sampleCount \oplus \{newUser? \mapsto 0\} \\
result! = success
\end{schema}

\subsubsection*{Error Case}

\begin{schema}{UserAlreadyExists}
\Xi NeuroScript \\
newUser? : USER \\
result! : REPORT
\where
newUser? \in registeredUsers \\
result! = alreadyExists
\end{schema}

\[
RRegisterUser \defs RegisterUser \lor UserAlreadyExists
\]

\subsection{User Login}

\begin{schema}{LoginUser}
\Xi NeuroScript \\
user? : USER \\
result! : REPORT
\where
user? \in registeredUsers \\
result! = success
\end{schema}

\begin{schema}{LoginFail}
\Xi NeuroScript \\
user? : USER \\
result! : REPORT
\where
user? \notin registeredUsers \\
result! = notFound
\end{schema}

\[
RLoginUser \defs LoginUser \lor LoginFail
\]

\subsection{Upload Handwriting Sample}

\begin{schema}{UploadSample}
\Delta NeuroScript \\
user? : USER \\
sid? : SAMPLEID \\
result! : REPORT
\where
user? \in registeredUsers \\
sid? \notin dom samples \\
sampleCount~user? < MAX\_SAMPLES\_PER\_USER \\
samples' = samples \cup \{sid? \mapsto user?\} \\
sampleStatus' = sampleStatus \cup \{sid? \mapsto Pending\} \\
features' = features \\
predictions' = predictions \\
explanations' = explanations \\
sampleCount' = sampleCount \oplus \{user? \mapsto sampleCount~user? + 1\} \\
result! = success
\end{schema}

\subsection{Preprocess and Extract Features}

\begin{schema}{ExtractFeatures}
\Delta NeuroScript \\
sid? : SAMPLEID \\
fv? : FEATUREVECTOR \\
result! : REPORT
\where
sid? \in dom samples \\
sampleStatus~sid? = Pending \\
features' = features \cup \{sid? \mapsto fv?\} \\
sampleStatus' = sampleStatus \oplus \{sid? \mapsto Processing\} \\
predictions' = predictions \\
explanations' = explanations \\
samples' = samples \\
registeredUsers' = registeredUsers \\
sampleCount' = sampleCount \\
result! = success
\end{schema}

\subsection{Classify Sample}

\begin{schema}{ClassifySample}
\Delta NeuroScript \\
sid? : SAMPLEID \\
cond? : CONDITION \\
result! : REPORT
\where
sid? \in dom features \\
sid? \notin dom predictions \\
predictions' = predictions \cup \{sid? \mapsto cond?\} \\
sampleStatus' = sampleStatus \oplus \{sid? \mapsto Analysed\} \\
features' = features \\
samples' = samples \\
registeredUsers' = registeredUsers \\
sampleCount' = sampleCount \\
explanations' = explanations \\
result! = success
\end{schema}

\subsection{Generate Explanation}

\begin{schema}{GenerateExplanation}
\Delta NeuroScript \\
sid? : SAMPLEID \\
exp? : EXPLANATION \\
result! : REPORT
\where
sid? \in dom predictions \\
sid? \notin dom explanations \\
explanations' = explanations \cup \{sid? \mapsto exp?\} \\
registeredUsers' = registeredUsers \\
sampleCount' = sampleCount \\
sampleStatus' = sampleStatus \\
samples' = samples \\
features' = features \\
predictions' = predictions \\
result! = success
\end{schema}

\subsection{Retrieve Prediction Result}

\begin{schema}{GetPrediction}
\Xi NeuroScript \\
sid? : SAMPLEID \\
condition! : CONDITION \\
result! : REPORT
\where
sid? \in dom predictions \\
condition! = predictions~sid? \\
result! = success
\end{schema}

\subsection{Search Samples by Condition}

\begin{schema}{SearchByCondition}
\Xi NeuroScript \\
user? : USER \\
cond? : CONDITION \\
results! : \power SAMPLEID \\
result! : REPORT
\where
user? \in registeredUsers \\
results! = \{ s : dom predictions | \\
\t1 samples~s = user? \land predictions~s = cond? \} \\
result! = success
\end{schema}

\subsection{Archive Sample}

\begin{schema}{ArchiveSample}
\Delta NeuroScript \\
sid? : SAMPLEID \\
result! : REPORT
\where
sid? \in dom samples \\
sampleStatus~sid? = Analysed \\
sampleStatus' = sampleStatus \oplus \{sid? \mapsto Archived\} \\
registeredUsers' = registeredUsers \\
samples' = samples \\
features' = features \\
predictions' = predictions \\
sampleCount' = sampleCount \\
explanations' = explanations \\
result! = success
\end{schema}

\subsection{Delete User Account}

\begin{schema}{DeleteUser}
\Delta NeuroScript \\
target? : USER \\
result! : REPORT
\where
target? \in registeredUsers \\
registeredUsers' = registeredUsers \setminus \{target?\} \\
samples' = \{ s : dom samples | samples~s \neq target? \} \ndres samples \\
sampleStatus' = \{ s : dom samples | samples~s \neq target? \} \ndres sampleStatus \\
features' = \{ s : dom samples | samples~s \neq target? \} \ndres features \\
predictions' = \{ s : dom samples | samples~s \neq target? \} \ndres predictions \\
explanations' = \{ s : dom samples | samples~s \neq target? \} \ndres explanations \\
result! = success
\end{schema}

\section{Summary of Schemas}

\begin{longtable}{|p{4cm}|p{3cm}|p{7cm}|}
\hline
\textbf{Schema} & \textbf{Type} & \textbf{Purpose} \\
\hline
NeuroScript & State Space & Core system state and invariants \\
\hline
InitNeuroScript & Initial State & All sets empty at start-up \\
\hline
RegisterUser & $\Delta$ Op & Add a new user with a role \\
\hline
LoginUser & $\Xi$ Op & Authenticate existing user \\
\hline
UploadSample & $\Delta$ Op & Upload handwriting sample \\
\hline
ExtractFeatures & $\Delta$ Op & Preprocess and extract feature vector \\
\hline
ClassifySample & $\Delta$ Op & LSTM-based neurological prediction \\
\hline
GenerateExplanation & $\Delta$ Op & Produce SHAP/LIME explanation \\
\hline
GetPrediction & $\Xi$ Op & Query prediction result \\
\hline
SearchByCondition & $\Xi$ Op & Search samples by predicted condition \\
\hline
ArchiveSample & $\Delta$ Op & Archive analysed sample \\
\hline
DeleteUser & $\Delta$ Op & Remove user and associated data \\
\hline
\end{longtable}

\vspace{0.5cm}

All major operations have corresponding error schemas and are combined into robust operations using the Z disjunction operator $(\lor)$, ensuring every possible input scenario produces a defined output via the REPORT free type.

\end{document}